\subsection{Permanencia de diversas propiedades de los morfismos por descenso
fielmente plano.}
\label{subsection:IV.2.7-fr}

\begin{proposition}[2.7.1]
\label{IV.2.7.1-fr}
Sean $f:X\to Y$ un $S$-morfismo de $S$-preesquemas, $g:S'\to S$ un
morfismo fielmente plano y cuasicompacto, $X'=X_{(S')}$, $Y'=Y_{(S')}$,
$f'=f_{(S')}:X'\to Y'$. Consideremos, para un morfismo, la propiedad de ser:
\begin{enumerate}
\item[\textup{(i)}] separado;
\item[\textup{(ii)}] cuasiseparado;
\item[\textup{(iii)}] localmente de tipo finito;
\item[\textup{(iv)}] localmente de presentación finita;
\item[\textup{(v)}] de tipo finito;
\item[\textup{(vi)}] de presentación finita;
\item[\textup{(vii)}] propio;
\item[\textup{(viii)}] un isomorfismo;
\item[\textup{(ix)}] un monomorfismo;
\item[\textup{(x)}] una inmersión abierta;
\item[\textup{(xi)}] una inmersión cuasicompacta;
\item[\textup{(xii)}] una inmersión cerrada;
\item[\textup{(xiii)}] afín;
\item[\textup{(xiv)}] cuasiafín;
\item[\textup{(xv)}] finito;
\item[\textup{(xvi)}] cuasifinito;
\item[\textup{(xvii)}] entero.
\end{enumerate}
Entonces, si $\boldsymbol{P}$ designa una de las propiedades precedentes, para que $f$
posea la propiedad $\boldsymbol{P}$, es necesario y suficiente que $f'$ la posea.
\end{proposition}

Se ha demostrado en los capítulos I, II y IV, \S~1, que, si $f$ posee una de
las propiedades $\boldsymbol{P}$ precedentes, lo mismo
ocurre con $f'$ (sin hipótesis sobre el morfismo $g:S'\to S$). Queda, pues,
demostrar el recíproco; como la proyección $Y'\to Y$ es

\oldpage[IV-2]{30}
un morfismo fielmente plano y cuasicompacto (2.2.13 y 1.1.2), es posible
limitarse al caso en que $S=Y$, $S'=Y'$ en virtud de (\textbf I, 3.3.11).

(i) Decir que $f$ es separado significa que el morfismo diagonal
$\Delta_f:X\to X\times_Y X$ es \emph{cerrado}; como
$\Delta_{f'}=(\Delta_f)_{(Y')}$ (\textbf I, 5.3.4), si $\Delta_{f'}$ es cerrado, lo
mismo ocurre con $\Delta_f$ en virtud de (2.6.2), por lo que $f$ es separado.

(ii) ya se ha demostrado bajo hipótesis más débiles (1.2.5).

(iii) y (iv): La cuestión es evidentemente local en $X$ y
$Y$ y, teniendo en cuenta (2.2.12), basta, pues, demostrar el
